\safesection{Esfuerzo axil}
$N_x$ es la resultante de fuerzas normal a la sección.

Convención de la asignatura:\claimref{CLM-003}
\[
N_x>0\Longleftrightarrow\text{tracción},
\qquad
N_x<0\Longleftrightarrow\text{compresión}.
\]

\safesection{Esfuerzo cortante}
En un problema plano:
\[
C_y
\]
o, en muchos libros,
\[
V.
\]
En este manual se utilizan indistintamente:
\[
V\equiv C_y.
\]
La convención utilizada en las diapositivas puede recordarse como:
\[
C_y>0:\ \text{levanta la cara izquierda}.
\]

\safesection{Momento flector}
En el plano $xy$:
\[
M_z.
\]
La convención del curso es:
\[
M_z>0:\ \text{comprime la parte superior}.
\]

\safesection{Momento torsor}
El momento alrededor de la directriz:
\[
M_x=T.
\]
Según el criterio mostrado en el material docente, $T>0$ corresponde al sentido antihorario observado sobre la cara frontal.

\safesection{Esfuerzo interno no es tensión}
Un esfuerzo interno es una resultante:
\[
N,\quad V,\quad M,\quad T.
\]
Una tensión describe cómo se distribuye esa resultante sobre el área:
\[
\sigma,\quad\tau.
\]
Por ejemplo:
\[
N=\int_A\sigma_{xx}\dd A.
\]

\flowbox{$\text{esfuerzo}\neq\text{tensión}$}

\begin{errorbox}
Decir ``la tensión axil vale $30\ \mathrm{kN}$'' es dimensionalmente incorrecto. $30\ \mathrm{kN}$ es una fuerza interna; una tensión debería expresarse, por ejemplo, en $\mathrm{MPa}$.
\end{errorbox}

% ============================================================
\safechapter{Equilibrio estático, equilibrio elástico e isostaticidad}
% ============================================================

\safesection{Equilibrio global}
Para una estructura plana:
\[
\sum F_x=0,\qquad
\sum F_y=0,\qquad
\sum M=0.
\]
Estas ecuaciones permiten calcular reacciones cuando la estructura es estáticamente determinada.

\safesection{Equilibrio de cualquier parte}
No basta con que el conjunto esté en equilibrio. \textbf{Cualquier porción} de la estructura obtenida mediante un corte también debe estarlo. Ésta es la base del método de las secciones.

\safesection{Estructura isostática}
\index{estructura!isostática}Una estructura es isostática cuando todas sus reacciones y esfuerzos pueden determinarse únicamente mediante equilibrio.

Para cargas mecánicas ordinarias:
\flowbox{los esfuerzos de una estructura isostática no dependen de $E$, $A$ o $I$ para poder ser obtenidos por equilibrio}

\safesection{Estructura hiperestática}
\index{estructura!hiperestática}En una estructura hiperestática existen incógnitas redundantes. Equilibrio por sí solo no basta.

Debemos añadir compatibilidad de deformaciones y ley constitutiva. Por ello los esfuerzos pueden depender de
\[
E,\qquad A,\qquad I,\qquad GJ,\ldots
\]

\begin{idea}
La diferencia esencial es:
\begin{center}
\begin{tabular}{c@{\hspace{2cm}}c}
\fbox{\parbox{0.30\linewidth}{\centering \textbf{Isostática}\\[2mm]equilibrio}}
&
\fbox{\parbox{0.38\linewidth}{\centering \textbf{Hiperestática}\\[2mm]equilibrio + compatibilidad + material}}
\end{tabular}
\end{center}
\end{idea}

% ============================================================
\safechapter{Método de las secciones}
% ============================================================

\index{método de las secciones}El método fundamental para obtener esfuerzos internos es:
\begin{enumerate}
  \item calcular las reacciones;
  \item realizar un corte en una posición genérica $x$;
  \item conservar uno de los dos lados;
  \item sustituir la parte eliminada por $N$, $V$, $M$;
  \item imponer equilibrio;
  \item repetir cuando cambie la ley de carga.
\end{enumerate}

